\begin{tikzpicture}
\matrix[tight matrix,nodes={style={text width=7cm,minimum height=.6cm,font=\small}}] (stack) {
    printf return address \\
    \myemph<4>{printf argument 7}/buffer start \& \tt "\myemph<2>{\%.419873}" \\
    \myemph<4>{printf argument 8} \& \tt "\myemph<2>{4u}\myemph<3>{\%c\%c\%c}" \\
    \myemph<4>{printf argument 9} \& \tt "\myemph<3>{\%c}\myemph<4>{\%c\%c\%c}" \\
    \myemph<4>{printf argument 10} \& \tt "\myemph<4>{\%c}\myemph<5>{\%ln}\myemph<6>{...}" \\
    \myemph<5>{printf argument 11} \& target \tt 0x7fffffffecf8 \\
    \ldots \& \ldots \\
};
\tikzset{
    note/.style={draw=red,very thick,anchor=north,align=left,at={([yshift=-.5cm]stack.south)}},
}
\begin{visibleenv}<2>
\node[note] {
    write unsigned number with 4198734 digits of percision \\
    result: \%rsi (printf arg 2) output \\
    padded to 4198734 digits with zeroes
};
\end{visibleenv}
\begin{visibleenv}<3>
\node[note] {
    one char (byte) based on printf args 3, 4, 5, 6 \\
    (\%rdx, \%rcx, \%r8, \%r9) \\
};
\end{visibleenv}
\begin{visibleenv}<4>
\node[note] {
    one char (byte) based on printf args 7, 8, 9, 10 \\
    (stack locations)
};
\end{visibleenv}
\begin{visibleenv}<5>
\node[note] {
    store number of bytes printed into printf arg 11 \\
    \texttt{l} indicates that it a long (not int) \\
    total bytes = 4198734 (\%u) + 8 (\%c $\times$ 8) = 0x401156
};
\end{visibleenv}
\begin{visibleenv}<6>
\node[note] {
    extra data just to ensure the target address \\
    is positioned correctly
};
\end{visibleenv}
\end{tikzpicture}